\section{Linear Octree Construction}

To build the Linear Octree, we start with the Morton code array and the particle ID array. Both arrays are already sorted according to the Morton codes. We use the following arrays to store the tree structure:

\begin{itemize}
    \item \textbf{node\_start} - Stores the starting particle index of the node.
    \item \textbf{node\_count} - Stores the number of particles in the node.
    \item \textbf{node\_parent} - Stores the parent node index.
    \item \textbf{child\_idx} - Stores the first child index of the node.
\end{itemize}

Initially, index 0 is the root node. It contains all the particles and has no parent. Therefore, we set:

\begin{lstlisting}[language=C]
node_start[0] = 0;
node_count[0] = num_particles;
node_parent[0] = -1;
\end{lstlisting}

We also keep track of the current level of the tree. The root is at level 0, so we initialize \textbf{level\_start\_idx = 0}, \textbf{num\_nodes\_in\_level = 1}, and \textbf{current\_level = 0}. A global counter on the GPU is used to keep track of the total number of nodes. We also use the \textbf{out\_level\_bounds} array to store the starting index of each level. The difference between two consecutive values gives the number of nodes in that level.

The tree is constructed level by level. The loop continues while the current level contains nodes and the maximum Morton-code depth has not been reached. The depth is limited to 21 levels because the Morton code uses three 21-bit coordinates, giving 63 useful bits.

For each level, we launch the CUDA kernel with 256 threads per block. Since the nodes are processed level by level, the global node index is calculated using the start index of the current level:

\begin{lstlisting}[language=C]
int idx = level_start_idx +
          (blockIdx.x * blockDim.x + threadIdx.x);
\end{lstlisting}

If the index is outside the current level, the thread returns.

For every valid node, we check its particle count. If the count is less than or equal to \textbf{N\_crit}, the node becomes a leaf and its \textbf{child\_idx} is set to -1. Otherwise, the node is subdivided into 8 children.

The child nodes need to be given unique positions in the global node arrays. Since several CUDA threads can try to create children at the same time, we use \textbf{atomicAdd} on the global node counter:

\begin{lstlisting}[language=C]
int first_child = atomicAdd(d_global_node_counter, 8);
\end{lstlisting}

This gives each parent a separate group of 8 continuous indices. The value returned by \textbf{atomicAdd} is used as the first child index.

\subsection{Handling the Morton Codes}

The Morton codes are stored as 64-bit unsigned integers. Three 21-bit coordinates are packed into the code, so 63 bits are used. Every tree level corresponds to 3 Morton-code bits, which gives a maximum depth of 21 levels.

To find the bits needed for the next level, we calculate:

\begin{lstlisting}[language=C]
int prefix_bits = (current_level + 1) * 3;
int shift_amount = 63 - prefix_bits;
\end{lstlisting}

The parent prefix is obtained by shifting the Morton code. We then shift this prefix left by 3 bits. These next three bits identify which of the 8 octants a particle belongs to.

\begin{lstlisting}[language=C]
uint64_t parent_prefix =
    morton[parent_start] >> (current_level * 3);

uint64_t child_base_prefix = parent_prefix << 3;
\end{lstlisting}

The Morton code therefore gives us a way to identify the child cells without explicitly storing the spatial coordinates of every particle during the tree construction.

\subsection{Assigning Particles to Children}

For every subdivided node, we process its 8 children in order. The particle range of each child is found from the sorted Morton codes.

For children 0 through 6, we form a target prefix for the next child and use \textbf{find\_lower\_bound()} to find the boundary in the sorted Morton-code array. For the last child, the end of the parent's particle range is used directly.

\begin{lstlisting}[language=C]
if (i == 7) {
    next_start = parent_start + parent_count;
} else {
    uint64_t target_prefix =
        child_base_prefix | (i + 1);

    next_start = find_lower_bound(
        morton,
        current_start,
        parent_start + parent_count,
        target_prefix,
        shift_amount);
}
\end{lstlisting}

The \textbf{find\_lower\_bound()} function uses a binary search on the sorted Morton codes. It finds the position where the next child prefix starts. This gives us the boundary between two child particle ranges.

After finding the boundary, we store the starting position and particle count of the child:

\begin{lstlisting}[language=C]
d_node_start[c_idx] = current_start;
d_node_count[c_idx] = next_start - current_start;

current_start = next_start;
\end{lstlisting}

The same process is repeated for all 8 children.

\subsection{Returning to the Host}

After the kernel finishes, the host code synchronizes the device and copies the updated global node count back to the CPU. This gives the total number of nodes created so far.

For the next level, \textbf{level\_start\_idx} is moved to the first node of that level, and \textbf{num\_nodes\_in\_level} is updated. The current level is then incremented. If the new level contains nodes, its starting index is added to the level-bounds array.

This process continues until there are no more nodes to process or the maximum Morton-code depth is reached.

At the end of this process, the tree structure is complete. We know the particle range of every node, its parent, and its children. The next step is to map this structure to its actual spatial geometry.
